% !TEX root = ../main.tex
% =====================================================================
% 2. BACKGROUND -- page budget 1.0
%
% Only what a reader needs to follow Sections 5-8: the MIL view, the
% baseline head we compare everything against, and the two metrics.
% The notation defined here is never redefined later.
% =====================================================================

\section{Background}
\label{sec:bg}

\subsection{Time series as bags of time steps}
\label{sec:bg:mil}

A univariate series $\bag = (\inst_1, \dots, \inst_{\nsteps})$ carries one label
$y \in \{1,\dots,\nclz\}$, but in most datasets only a few time steps hold the evidence for it. That
is exactly the structure of Multiple Instance Learning \citep{dietterich1997mil}, where a labelled
\emph{bag} contains unlabelled \emph{instances}: take the series as the bag and each time step as an
instance, and recovering the explanation becomes the standard MIL problem of recovering
instance-level contributions from a bag-level label.

An encoder maps every step to an embedding, $\emb_j = f_\theta(\bag)_j \in \R^{\dmodel}$, and a MIL
pooling head reduces the $\nsteps$ embeddings to a bag prediction $\bagpred \in \R^{\nclz}$.
Attention-based MIL \citep{ilse2018attention} learns how much each instance counts, and because the
attention weight is the quantity the prediction is built from, it can be read directly as an
importance score. The most active application of this idea is histopathology, where a slide is a bag
and a patch is an instance; two methods from that literature -- multi-scale fusion
\citep{hashimoto2020msdamil} and the dual-stream critical-instance mechanism of
\citet{li2021dsmil} -- are adapted to time series in \Cref{sec:design}.

\subsection{The baseline: conjunctive pooling}
\label{sec:bg:conj}

\millet{} \citep{early2024millet} pairs an InceptionTime encoder \citep{fawaz2020inceptiontime} with
\emph{conjunctive pooling}, which combines one attention gate with a per-time-step classifier:
\begin{equation}
  a_j = \sig\!\big(\mathbf{w}_a^\top \tanh(W_a \emb_j)\big) \in [0,1],
  \qquad
  \instpred_j = W_c \emb_j + \mathbf{b}_c \in \R^{\nclz},
  \label{eq:conj_gate}
\end{equation}
\begin{equation}
  \bagpred^{\,k} = \frac{1}{\nsteps}\sum_{j=1}^{\nsteps} a_j \instpred_j^{\,k},
  \qquad
  \interp_{k,j} = a_j \instpred_j^{\,k} .
  \label{eq:conj_bag}
\end{equation}
The interpretation $\interp_{k,j}$ is made of the same terms that are summed to give
$\bagpred^{\,k}$, so the explanation cannot disagree with the prediction: it \emph{is} the
prediction, decomposed. Three properties of \Cref{eq:conj_bag} bound what the head can express, and
each one motivates a family of heads in \Cref{sec:design}. \textbf{L1, one gate for all classes:}
$a_j$ is a single scalar per step, so a step cannot support one class while arguing against another.
\textbf{L2, weights not normalised over time:} the gate is a sigmoid and the sum is divided by
$\nsteps$, so the magnitude of $\bagpred^{\,k}$ depends on the series length -- awkward on an archive
whose lengths span $15$ to $2{,}844$. \textbf{L3, a fixed mean:} evidence concentrated in a few steps
of a long series contributes a vanishing share of the average, and a fixed operator cannot adapt to
datasets where the evidence is instead spread out.

\subsection{Measuring explanation quality}
\label{sec:bg:metrics}

We use the two metrics of \citet{early2024millet}, computed with their implementation, because our
purpose is to compare against their method on their own terms.

\paragraph{AOPCR (faithfulness to the model).}
Steps are ranked by the model's own interpretation, masked one at a time in that order, and the
confidence in the original prediction is tracked; a faithful ranking makes it fall fast. Writing
$O^k$ for the induced ranking and $\bag^{O^k}_j$ for the series with its $j$ most important steps
masked \citep{samek2017aopc},
\begin{equation}
  \mathrm{AOPC}(\bag, O^k) = \frac{1}{\nsteps-1}\sum_{j=1}^{\nsteps-1}
  \Big[\bagpred^{k}(\bag) - \bagpred^{k}\big(\bag^{O^k}_{j}\big)\Big],
  \label{eq:aopc}
\end{equation}
and AOPCR reports the margin of \Cref{eq:aopc} over three random maskings, so that a merely fragile
model is not rewarded. AOPCR needs no ground truth and is therefore available on every dataset. It
is also, as \Cref{sec:aopcr} shows, \textbf{unnormalised} -- it is built from raw confidence
differences, so its scale follows the scale of the logits.

\paragraph{\ndcg{} (correctness against ground truth).}
Normalised Discounted Cumulative Gain \citep{jarvelin2002ndcg} asks whether the steps that
\emph{genuinely} matter appear at the top of the ranking. It needs per-time-step ground truth, which
only a synthetic generator provides, but it is normalised to $[0,1]$ and is therefore the metric we
lean on whenever the two disagree.
